\documentclass[a4,11pt]{aleph-notas}
% Se puede ver la documentación aquí:
% https://github.com/alephsub0/LaTeX_aleph-notas

% -- Paquetes adicionales
\usepackage{enumitem}
\usepackage{url}
\usepackage{array}
\usepackage{booktabs}
\usepackage{longtable}
\usepackage{aleph-comandos}
\hypersetup{
    urlcolor=blue,
    linkcolor=blue,
}


% -- Datos
\institucion{Facultad de Ciencias Exactas, Naturales y Ambientales}
\carrera{Catálogo STEM}
\asignatura{Matemática III}
\tema{Plan tema 16: Integrales de superficie}
\autor{Andrés Merino}
\fecha{Periodo 2026-1}

\logouno[0.16\textwidth]{Logos/logoPUCE_04_ac}
\definecolor{colortext}{HTML}{1957d1}
\definecolor{colordef}{HTML}{1957d1}
\fuente{montserrat}


\begin{document}

\encabezado

%%%%%%%%%%%%%%%%%%%%%%%%%%%%%%%%%%%%%%%%
\section*{Resultado de Aprendizaje}
%%%%%%%%%%%%%%%%%%%%%%%%%%%%%%%%%%%%%%%%

%%%%%%%%%%%%%%%%%%%%%%%%%%%%%%%%%%%%%%%%
\subsection*{RdA de la asignatura:}
\begin{itemize}[leftmargin=*]
    \item \textbf{RdA 1:} Identificar conceptos, teoremas y propiedades de funciones vectoriales, derivadas parciales e integral vectorial.
    \item \textbf{RdA 2:} Calcular derivadas parciales e integrales con asistentes computacionales.
    \item \textbf{RdA 3:} Aplicar Cálculo Vectorial para resolver problemas reales.
\end{itemize}

%%%%%%%%%%%%%%%%%%%%%%%%%%%%%%%%%%%%%%%%
\subsection*{RdA del tema:}
\begin{itemize}[leftmargin=*]
    \item Parametrizar superficies en el espacio y calcular sus vectores tangentes y normal.
    \item Calcular el área de una superficie suave mediante la norma del vector normal.
    \item Calcular integrales de superficie de campos escalares sobre superficies parametrizadas.
    \item Calcular integrales de superficie de campos vectoriales (flujo) y reconocer el efecto de la orientación.
\end{itemize}

%%%%%%%%%%%%%%%%%%%%%%%%%%%%%%%%%%%%%%%%
\section*{Introducción}
%%%%%%%%%%%%%%%%%%%%%%%%%%%%%%%%%%%%%%%%

\paragraph{Pregunta inicial:}
Imaginemos el agua de un río que atraviesa una red colocada en diagonal a la corriente. La cantidad de agua que pasa por la red depende de su forma, de su inclinación y de la rapidez del agua en cada punto. ¿Cómo podríamos medir la cantidad total de agua que atraviesa una superficie curva?

%%%%%%%%%%%%%%%%%%%%%%%%%%%%%%%%%%%%%%%%
\section*{Desarrollo}
%%%%%%%%%%%%%%%%%%%%%%%%%%%%%%%%%%%%%%%%

%%%%%%%%%%%%%%%%%%%%%%%%%%%%%%%%%%%%%%%%
\subsection*{Actividad 1: Superficies parametrizadas e integrales de campos escalares}
Se introduce la parametrización de superficies, los vectores tangentes y normal y el cálculo de área, para luego definir la integral de superficie de un campo escalar.

\paragraph{¿Cómo lo haremos?}
\begin{itemize}[leftmargin=*]
    \item \textbf{Motivación:} mostrar cómo una superficie en el espacio (un plano inclinado, un cilindro, una esfera) se describe mediante dos parámetros, de forma análoga a como una curva se describe con uno.
    \item \textbf{Clase magistral:} se definen parametrización, curvas coordenadas, vectores tangentes $T^1$, $T^2$ y el vector normal $N=T^1\times T^2$; se obtiene el área como $\iint_\Omega \|N(u,v)\|\,dudv$; se define la integral de superficie de campos escalares con la fórmula $\iint_\Sigma f\,d\alpha=\iint_\Omega f(\alpha(u,v))\|N(u,v)\|\,dudv$. Se usa el resumen \href{https://andres-merino.github.io/MatematicaIII-05-N0051/01Resumen/Resumen16.pdf}{Resumen16.pdf}.
    \item \textbf{Implementación computacional:} visualizar superficies parametrizadas y calcular áreas e integrales de superficie con un asistente computacional.
    \item \textbf{Resolución de ejercicios:} calcular áreas e integrales de superficie de campos escalares usando \href{https://andres-merino.github.io/MatematicaIII-05-N0051/04Ejercicios/Ejercicios16.pdf}{Ejercicios16.pdf}.
\end{itemize}

\paragraph{Verificación de aprendizaje:}
\begin{itemize}[leftmargin=*]
    \item Para la parametrización $\alpha(u,v)=(u,v,u^2+v^2)$, calcule los vectores tangentes y el vector normal en un punto.
    \item Calcule el área de la porción del plano $z=x+y$ sobre el cuadrado $[0,1]\times[0,1]$.
    \item ¿Qué representa geométricamente la norma del vector normal $\|N(u,v)\|$?
\end{itemize}

%%%%%%%%%%%%%%%%%%%%%%%%%%%%%%%%%%%%%%%%
\subsection*{Actividad 2: Integrales de superficie de campos vectoriales}
Se define la integral de superficie de un campo vectorial (flujo) y se analiza el papel de la orientación de la superficie.

\paragraph{¿Cómo lo haremos?}
\begin{itemize}[leftmargin=*]
    \item \textbf{Motivación:} retomar la pregunta inicial e interpretar la integral de superficie de un campo vectorial como el flujo de un fluido a través de la superficie.
    \item \textbf{Clase magistral:} se define la integral de superficie de campos vectoriales con la fórmula $\iint_\Sigma F\cdot d\alpha=\iint_\Omega F(\alpha(u,v))\cdot N(u,v)\,dudv$; se discuten las parametrizaciones equivalentes y el efecto de invertir la orientación. Se usa el resumen \href{https://andres-merino.github.io/MatematicaIII-05-N0051/01Resumen/Resumen16.pdf}{Resumen16.pdf}.
    \item \textbf{Resolución de ejercicios:} calcular el flujo de campos vectoriales a través de superficies usando \href{https://andres-merino.github.io/MatematicaIII-05-N0051/04Ejercicios/Ejercicios16.pdf}{Ejercicios16.pdf}.
\end{itemize}

\paragraph{Verificación de aprendizaje:}
\begin{itemize}[leftmargin=*]
    \item Calcule el flujo del campo $F(x,y,z)=(0,0,1)$ a través de la porción del plano $z=0$ sobre $[0,1]\times[0,1]$.
    \item ¿Cómo cambia el valor del flujo si se invierte la orientación de la superficie?
    \item Explique la diferencia entre la integral de superficie de un campo escalar y la de un campo vectorial.
\end{itemize}

%%%%%%%%%%%%%%%%%%%%%%%%%%%%%%%%%%%%%%%%
\section*{Cierre}
%%%%%%%%%%%%%%%%%%%%%%%%%%%%%%%%%%%%%%%%

\paragraph{Tarea:}
Resolver del libro \href{https://catalogobiblioteca.puce.edu.ec/cgi-bin/koha/opac-detail.pl?biblionumber=83405&query_desc=su%3A%22An%C3%A1lisis%20vectorial%22}{Cálculo Vectorial de Colley}: sección 6.1, ejercicios 15--25 (impares).

\paragraph{Pregunta de investigación:}
\begin{enumerate}[leftmargin=*]
    \item ¿Qué es una superficie no orientable? Investigar el ejemplo de la banda de Möbius.
    \item ¿En qué aplicaciones físicas aparece el flujo a través de una superficie? Investigar un ejemplo en electromagnetismo o en dinámica de fluidos.
\end{enumerate}

\paragraph{Para el próximo tema:}
Ver los videos:
\begin{itemize}[leftmargin=*]
    \item \href{https://youtu.be/a4cFvAAt9zU?si=Ot8O8wnYy1cWk0lN}{Green's theorem proof}.
\end{itemize}


\end{document}
